\documentclass[11pt]{article}

\usepackage[margin=1in]{geometry}
\usepackage{amsmath,amssymb,amsthm,mathtools}
\usepackage{tikz-cd}
\usepackage{hyperref}
\usepackage{microtype}

\newtheorem{theorem}{Theorem}[section]
\newtheorem{corollary}[theorem]{Corollary}
\newtheorem{proposition}[theorem]{Proposition}
\newtheorem{lemma}[theorem]{Lemma}
\newtheorem{definition}[theorem]{Definition}
\newtheorem{remark}[theorem]{Remark}

\newcommand{\N}{\mathbb N}
\newcommand{\Q}{\mathbb Q}
\newcommand{\Z}{\mathbb Z}
\newcommand{\im}{\operatorname{im}}
\newcommand{\Spec}{\operatorname{Spec}}
\newcommand{\Supp}{\operatorname{Supp}}
\newcommand{\Real}{\operatorname{Real}}

\title{Fibrewise Exactification and Reciprocal Interval Profiles\\
\large A categorical proof of Erd\H{o}s Problem 289}
\author{Yuren Tang}
\date{4 September 2026}

\begin{document}
\maketitle

\begin{abstract}
We prove Erd\H{o}s Problem 289.  We work in the stronger subsystem in which every connected component has size $2$ or $3$.  If $b_2$ and $b_3$ denote the numbers of binary and ternary components, respectively, then there exists $d_*$ such that for every fixed $d\ge d_*$, all sufficiently large binary counts $b$ occur in an exact reciprocal representation of $1$ with component profile $(b,d)$.  Consequently every sufficiently large total component count occurs.

The proof is organized around a regular-categorical exactification mechanism.  Correction requirements are pullbacks, physical coverage is expressed by regular epimorphisms from realizer pullbacks, and physical composition is performed on literal matched-realizer objects before passage to regular-image support.  The target $1$ is centered by the universal quotient $\Q\to\Q/\Z$.  The application-specific root consists of two filtered responses: a filtered affine ray and a late-source targeted binary-residue response.  Their source filters are matched by a $\operatorname{Tendsto}$ pullback, after which exactification is the direct cancellation of opposite cosets in $\Q/\Z$.  The small-resource quantifier necessarily precedes the eventual source quantifier; fixed-source arbitrarily-small positive saturation is false.
\end{abstract}

\section{Introduction}

For a finite set $S\subset\N_{>0}$ write
\[
W(S)=\sum_{n\in S}\frac1n.
\]
Erd\H{o}s Problem 289 asks whether, for every sufficiently large integer $k$, one can find exactly $k$ pairwise disjoint and non-adjacent finite intervals of positive integers, each of length at least $2$, whose reciprocal weights sum to $1$.

Rather than treating an interval decomposition as primitive data, we regard a finite subset of the path graph on $\N_{>0}$ as the physical state and obtain its blocks from the connected-component quotient.  We restrict to the subsystem in which every component has cardinality $2$ or $3$.

For such a state let
\[
b_2(S)=\#\{\text{components of size }2\},
\qquad
b_3(S)=\#\{\text{components of size }3\}.
\]
The intrinsic additive exact label is the component profile
\[
\pi(S)=(b_2(S),b_3(S))\in\N^2.
\]
The usual number of intervals is only the consumer
\[
g(S)=b_2(S)+b_3(S).
\]

\begin{theorem}[Fibrewise profile saturation]\label{thm:fibrewise}
There exists $d_*\in\N$ such that for every $d\ge d_*$ there is $b_0(d)$ with the following property.  For every $b\ge b_0(d)$ there exists a finite support $S\subset\N_{>0}$ such that
\begin{enumerate}
\item every connected component of $S$ has size $2$ or $3$;
\item distinct components are non-adjacent;
\item $W(S)=1$;
\item $(b_2(S),b_3(S))=(b,d)$.
\end{enumerate}
\end{theorem}

\begin{corollary}[Erd\H{o}s Problem 289]\label{cor:erdos289}
For every sufficiently large integer $k$, there exist exactly $k$ pairwise disjoint and non-adjacent finite intervals of positive integers, each of length at least $2$, whose reciprocal weights sum to $1$.
\end{corollary}

\begin{proof}
Fix any $d\ge d_*$.  By Theorem~\ref{thm:fibrewise}, every sufficiently large $b$ occurs with profile $(b,d)$.  The total number of components is $b+d$, hence every sufficiently large total component count occurs.  The components already have sizes $2$ or $3$.
\end{proof}

For every finite $D\ge d_*$, the theorem also gives a simultaneous finite-defect strengthening: all sufficiently large total grades can be realized at every prescribed ternary count $d\in[d_*,D]$, with the configuration allowed to depend on $d$.

\subsection{The proof architecture}

The shortest active dependency graph is
\[
\begin{tikzcd}[column sep=large,row sep=large]
& \text{regular correction / target-quotient calculus} \arrow[d] & \\
& \text{reciprocal component-profile specialization} \arrow[dl] \arrow[dr] & \\
\text{filtered affine ray} \arrow[dr]
&& \text{late-source targeted residue response} \arrow[dl] \\
& \text{source-filter pullback and intersection} \arrow[d] & \\
& \text{defectwise coset cancellation} \arrow[d] & \\
& \{b:(b,d)\text{ is exact}\}\text{ is final} \arrow[d] & \\
& \text{projection to total component count.} &
\end{tikzcd}
\]

The root retains neither a common tail across different defects nor a common affine correction, affine grade, or endpoint.  What it does retain is one cofinal affine source germ, because arbitrarily small positive tail mass is obtained by moving the source late.

\subsection{Why the source must move}

Let
\[
H_n=n^{-1}\Z/\Z\le\Q/\Z.
\]
If $0<W<1$ and the source residue of $W$ modulo $H_n$ is zero, then $[W]\in H_n$ and therefore
\[
W\in\left\{\frac1n,\frac2n,\ldots,\frac{n-1}{n}\right\}.
\]
Hence
\[
W\ge\frac1n.
\]
Thus no positive-grade response from a fixed source $H_n$ can have zero source residue and arbitrarily small positive mass.  The correct saturation theorem consequently has the quantifier order
\[
\varepsilon>0
\quad\text{before}\quad
\forall^\infty\text{ source}.
\]
This simple mesh obstruction is the reason a source filter survives in the final architecture.

\subsection{Compose before reflect}

Physical compatibility belongs to actual witnesses, not to aggregate labels.  We retain literal witness tuples until physical composition has been certified and only afterwards pass to regular-image support:
\[
\boxed{
\text{witness tuples}
\longrightarrow
\text{physical composition}
\longrightarrow
\text{regular-image support}
\longrightarrow
\text{profile calculus}.}
\]

\section{Physical states and component profiles}

Let $P$ be the path graph on $\N_{>0}$.  For a finite support $S$, let $E(S)$ be the edge object of the induced subgraph.  Its connected-component object is the coequalizer
\[
E(S)\rightrightarrows S\longrightarrow\pi_0(S).
\]
Admissibility requires every component to have cardinality in $\{2,3\}$ and distinct components to be separated by at least one unused vertex.

The reciprocal weight and component profile are additive on physically compatible unions:
\[
W(S)=\sum_{n\in S}\frac1n,
\qquad
\pi(S)=(b_2(S),b_3(S)).
\]
Write $C^{23}$ for this partial physical state system and
\[
\Sigma^{23}_{1}
=
\im\bigl(C^{23}_{W=1}\xrightarrow{\pi}\N^2\bigr)
\]
for its exact component-profile spectrum.

The ordinary grade map factors through the consumer triangle
\[
\begin{tikzcd}
C^{23}_{W=1} \arrow[r,"{(b_2,b_3)}"] \arrow[dr,"g"']
& \N^2 \arrow[d,"+"]\\
& \N.
\end{tikzcd}
\]
Consequently the ordinary exact-grade spectrum is the regular-image projection of $\Sigma^{23}_{1}$.

\section{The regular correction engine}

We record only the generic structure subsequently used by the local-to-global provider proof.

\subsection{Required objects and physical covers}

Work in a regular category $\mathcal E$.  Let $V$ be an exact-value carrier, $L$ an internal monoid of exact auxiliary labels, and
\[
W:C\to V,
\qquad
\lambda:C\to L
\]
be physical observations.

For one internal group homomorphism $u:Q_0\to Q_1$, a parameterized correction over a base $S$ consists of maps
\[
d:S\to UQ_1,
\qquad
\ell:S\to L.
\]
Put $\widetilde Q_i=UQ_i\times L$.  Its required source object is the pullback
\[
\begin{tikzcd}
T_S(\phi) \arrow[r] \arrow[d]
& \widetilde Q_0 \arrow[d,"Uu\times\mathrm{id}"]\\
S \arrow[r,"{(d,\ell)}"']
& \widetilde Q_1.
\end{tikzcd}
\]

For a physical family $F\to S$, its realizer object is the corresponding pullback against the physical observation.  Coverage is exactly
\[
\Real_F(\phi)\twoheadrightarrow T_S(\phi)
\]
regular epimorphic.

\subsection{Literal matching and physical completion}

For composable corrections $\phi,\psi$ with composite $\theta$, a point of $T_S(\theta)$ canonically forces the second requirement.  Pull a second realizer back along this forced map; the resulting actual second state then canonically forces the first requirement.  Pulling the first realizer cover back along that map yields
\[
\operatorname{Cmp}(F,R;\phi,\psi)
\twoheadrightarrow T_S(\theta).
\]

The generic theorem does not assume a globally defined multiplication on $C$.  A physical completion certificate is a requirement-preserving relation over the literal match object whose projection to $T_S(\theta)$ is regular epi and whose state map has the composite observation.  Compatible union in the reciprocal application supplies such certificates.

\subsection{Exact transfer}

For a target $\tau:S\to V$, exact states are the pullback
\[
\begin{tikzcd}
C_\tau \arrow[r] \arrow[d] & C \arrow[d,"W"]\\
S \arrow[r,"\tau"'] & V.
\end{tikzcd}
\]
The exact auxiliary-label spectrum is the regular image
\[
\Sigma_\tau=\im(C_\tau\to S\times L).
\]
A target correction has a canonical target section of its required object.  Pulling a correction cover back along this section and factoring the physical state map through $C_\tau$ forces the target label graph to factor through $\Sigma_\tau$.

\section{Witness semantics and support enrichment}

The exact component-profile label used by exact transfer is not identified with the ordered resource bound used inside the saturation proof.  For Erd\H{o}s Problem 289, Engine A uses
\[
L=P_{23}=\N^2.
\]
The centered saturation proof may use a larger realization monoid carrying binary count and a positive mass bound; its resource coordinate is forgotten after smallness has been consumed.

For a realization-label monoid $R$, labelled witness families lie in $\mathcal E/R$.  Their support is the regular image
\[
X\twoheadrightarrow\Supp(X)\hookrightarrow R.
\]
Witness convolution is formed before reflection, and subobject convolution is the regular image of label addition.  For a path $\pi$, a physical composition certificate is a span
\[
\begin{tikzcd}
& C_\pi \arrow[dl,two heads] \arrow[dr] &\\
T_\pi && W_{\operatorname{comp}\pi}
\end{tikzcd}
\]
over the realization label.  Therefore
\[
\Supp(T_\pi)\le\Supp(W_{\operatorname{comp}\pi}).
\]
Joining over paths gives the free enriched realization inequality used below the late-source Tail response.

\section{Target quotient and compact resolution}

For reciprocal target $1$, the marked map
\[
\Z\xrightarrow{1\mapsto1}\Q
\]
has universal target-annihilating quotient
\[
A:=\Q/\Z.
\]
The compact observation stages are the finitely generated subgroups of $A$, equivalently
\[
H_n=\ker([n]:A\to A).
\]
Their order is divisibility; gcd and lcm are coordinate descriptions of meet and join.

For $H\le A$ write $q_H:A\to A/H$.  If $r\in A/H$, the least compact extension killing $r$ is the pullback
\[
H[r]=q_H^{-1}(\langle r\rangle).
\]
This least absorber is used in the proof that the stronger historical BinaryTail theorem implies the targeted response below.  It is not required in the final coset-cancellation diagram.

Simple compact-stage factors are finite additive groups of prime order.  The local restricted-fold theorem therefore needs no chosen $\mathbf F_p$ basis.

\section{The late-source targeted response}

Fix a normalized raw-row system $R$.  Let
\[
S_R=\operatorname{EligibleBase}(R),
\qquad
\mathcal F_R=\operatorname{EligibleBaseFilter}(R)
\]
be its proper source filter.  For $\beta\in S_R$, let
\[
H_\beta\le A
\]
be the corresponding compact source stage.

For a row-transparent obstacle $O$, residue $r\in A/H_\beta$, tolerance $\varepsilon>0$, and binary count $h$, let
\[
\operatorname{LTarResp}_R(\varepsilon,c,\beta,O,r,h)
\]
be the normalized solution object of binary responses satisfying
\[
(b_2,b_3)=(h,0),
\qquad
W_+\le\varepsilon,
\qquad
q_{H_\beta}([W])=r,
\]
and the prescribed physical obstacle.

\begin{theorem}[Late-source targeted residue saturation]\label{thm:ltar}
For every $R$, every $\varepsilon>0$, and every transparency parameter $c>0$,
\[
\forall^\infty\beta\ [\mathcal F_R],
\quad
\forall O,
\quad
\forall r\in A/H_\beta,
\quad
\forall^\infty h,
\]
the projection
\[
\operatorname{LTarResp}_R(\varepsilon,c,\beta,O,r,h)\longrightarrow1
\]
is regular epimorphic.
\end{theorem}

The residue is quantified only after the late source is fixed, and the grade threshold may depend on the residue and obstacle.

\subsection{Derivation from the stronger BinaryTail theorem}

The historical BinaryTail theorem has quantifier order
\[
\forall R,\varepsilon,c,
\quad
\forall^\infty\beta[\mathcal F_R],
\quad
\forall O,
\quad
\forall K\in\operatorname{CompactRequest}(\beta),
\quad
\forall^\infty h.
\]
Fix a generalized late $\beta$ and $r\in A/H_\beta$.  Form the least absorber
\[
K_r=H_\beta[r].
\]
BinaryTail at $K_r$ supplies a full zero-correction realizer cover ending at some $G\ge K_r$.  Since $r$ is then a point of the relative kernel $G/H_\beta$, pulling the realizer regular epi back along $r$ gives the targeted residue response.  Regular epis are pullback-stable, so Theorem~\ref{thm:ltar} follows.

The arithmetic proof of BinaryTail is isolated later.  It is here that comparable-prime supply, signed-inverse rows, packing, the prime-order restricted-fold theorem, filtered torsor lifting, grade saturation, and resource decay occur.

\section{The filtered affine ray}

The affine provider must retain exactly enough information for its source stages to meet the good-source locus of Theorem~\ref{thm:ltar}.

\begin{definition}[Filtered affine ray]
A filtered affine ray consists of a proper filtered base $(B,\mathcal F_B)$, a source map
\[
\beta:B\to S_R
\]
satisfying
\[
\operatorname{Tendsto}(\beta,\mathcal F_B,\mathcal F_R),
\]
uniform constants $c>0$, $\eta>0$, and an eventually defined threshold $d_0:B\to\N$ such that, eventually in $b\in B$, every $d\ge d_0(b)$ admits a finite homogeneous physical family $F_{b,d}$ at source $H_b:=H_{\beta(b)}$ with
\[
\pi(F_{b,d})=(m_{b,d},d),
\]
for some $m_{b,d}$, whose centered-value image covers one affine coset
\[
q_{H_b}^{-1}(a_{b,d})\subseteq A
\]
regular epimorphically, whose branch values satisfy
\[
0<W_+<2-\eta,
\]
and whose complete physical footprint is row-transparent with constant $c$.
\end{definition}

\begin{theorem}[Affine-ray response]\label{thm:affray}
Such a filtered affine ray exists for the Erd\H{o}s Problem 289 physical system.
\end{theorem}

The theorem is a consumer-image projection of the stronger historical Prefix/MasterSlab construction.  Prefix supplies the proper base and the $\operatorname{Tendsto}$ source map.  MasterSlab supplies uniform $c,\eta$ and, eventually in the base, a parameterized affine cover over every finite defect interval.  For fixed $d$, set the slab ceiling equal to $d$, pull the parameterized cover back along the defect fibre $\{d\}$, and discard the other defect branches.  Common correction, common grade, common tail, selector presentation, and unused slab coordinates are not retained.

\section{Meeting the two source germs}

Work over the nonempty budget cover
\[
E_\eta=\{\varepsilon\in\Q_{>0}:0<\varepsilon<\eta\}\twoheadrightarrow1.
\]
For a generalized $\varepsilon$, Theorem~\ref{thm:ltar} supplies a good-source locus
\[
U_{\varepsilon,c}\in\mathcal F_R.
\]
The $\operatorname{Tendsto}$ law gives
\[
\beta^{-1}(U_{\varepsilon,c})\in\mathcal F_B.
\]
Intersect this with the eventual affine-ray locus.  Properness of $\mathcal F_B$ makes the intersection inhabited.  Work over the resulting generalized base point $b$ and put
\[
d_*:=d_0(b).
\]
No distinguished global source is selected or exported.

\section{Coset-cancellation exactification}

The final exactification is thinner than a full correction-cover composition theorem.

Fix one source subgroup $H\le A$ and one affine quotient point
\[
a\in A/H.
\]
The affine fibre is the pullback
\[
C_a=A\times_{A/H}1_a=q_H^{-1}(a).
\]
Suppose an affine family $F$ covers this fibre:
\[
R_F\twoheadrightarrow C_a.
\]
Let a targeted binary family $T$ satisfy
\[
q_H([W(T)])=-a.
\]
Then inversion gives
\[
T\longrightarrow C_a,
\qquad
t\longmapsto-[W(t)].
\]
Pull back the affine realizer cover:
\[
\begin{tikzcd}
M \arrow[r] \arrow[d,two heads]
& R_F \arrow[d,two heads]\\
T \arrow[r,"-{[W]}"']
& C_a.
\end{tikzcd}
\]
Every matched pair $(f,t)$ therefore satisfies
\[
[W(f)+W(t)]=0\in\Q/\Z.
\]
A physical completion certificate supplies an actual compatible union $f\oplus t$.

If
\[
0<W(f)<2-\eta,
\qquad
0\le W(t)<\varepsilon<\eta,
\]
then
\[
0<W(f\oplus t)<2.
\]
Since its class in $\Q/\Z$ is zero,
\[
W(f\oplus t)\in\Z\cap(0,2)=\{1\}.
\]
Thus
\[
\boxed{W(f\oplus t)=1.}
\]
No endpoint or least absorber is consumed in this final diagram.

\section{Proof of fibrewise profile saturation}

Fix the generalized base $b$ obtained by the source-filter intersection and let $d_* = d_0(b)$.  Fix any $d\ge d_*$.  The affine ray supplies
\[
(F_{b,d},a_{b,d},m_{b,d})
\]
at source $H_b$.

Apply Theorem~\ref{thm:ltar} at the same source and the complete affine obstacle, with target residue
\[
r=-a_{b,d}.
\]
For every sufficiently large binary tail count $h$, obtain a compatible binary response $t$ satisfying
\[
q_{H_b}([W(t)])=-a_{b,d},
\qquad
W_+(t)<\varepsilon,
\qquad
\pi(t)=(h,0).
\]
Coset cancellation gives a matched affine branch $f$ with exact reciprocal sum
\[
W(f\oplus t)=1.
\]
Component profiles add:
\[
\pi(f\oplus t)
=(m_{b,d},d)+(h,0)
=(m_{b,d}+h,d).
\]
Hence the exact binary-count fibre
\[
B_d=\{n:(n,d)\in\Sigma^{23}_{1}\}
\]
contains a translate of a final set and is therefore final.  This proves Theorem~\ref{thm:fibrewise}.

\section{Finite-defect uniformity}

For fixed $d\ge d_*$, let
\[
G_d=d+B_d
\]
be the set of total component counts realized at ternary count $d$.  Each $G_d$ is final.  Hence for every finite $D\ge d_*$,
\[
\bigcap_{d=d_*}^{D}G_d
\]
is final.  Thus every sufficiently large total component count is realizable at every prescribed $d\in[d_*,D]$.

\section{Proof boundary for late-source saturation}

The current proof of Theorem~\ref{thm:ltar} factors through the stronger BinaryTail theorem.  Below that theorem the ordinary mathematical inputs are:
\begin{enumerate}
\item comparable-prime supply;
\item transverse signed-inverse reciprocal supply;
\item bounded-degree quota packing;
\item restricted-fold surjectivity in a finite additive group of prime order;
\item filtered torsor lifting across non-split compact-stage extensions;
\item bounded-thinning/D-supercritical transference, interval overlap, and summable resource control.
\end{enumerate}
Historical rates such as $Q/\log Q$, $Q/p$, $Q^{-2}$, and $(\log Q)/Q^2$ occur only as proof certificates for these provider properties and do not survive the late-source response boundary.

A simple associated-graded coefficient does not canonically determine its lower shadow.  Thus a literal upper-triangular coefficient matrix would require a noncanonical splitting.  The present proof instead uses full local epimorphism and filtered torsor lifting.  Weakening this provider to target-induced matched images would require new extension-sensitive mathematics and is not part of the proof here.

\section{Further categorical compression}

For one specified target section, full correction-fibre coverage is generically stronger than necessary.  A sectionwise matched-image theorem requires only the forced requirements to land in the corresponding realizable-requirement images, together with a physical completion certificate.  This generic theorem is valid and conceptually useful.

The present reciprocal provider is deliberately used through its already established stronger full simple-fibre coverage.  No claim that the restricted-fold theorem can be removed is required for the proof.

\section{Provenance and formalization status}

The mathematical route developed independently before the proof claim posted on 4 September 2026.  A standalone concrete proof manuscript had been completed in mid-August, followed by categorical reorganization and an ongoing Lean formalization programme.  The present manuscript uses the later categorical reorganization rather than reproducing the older concrete proof architecture.

The Lean formalization is incomplete and is not claimed here as a formal verification of the final theorem.  Its role has nevertheless been useful in exposing witness-lifetime, compatibility, quantifier-order, and dependency issues that motivated several categorical normalizations and the late-source correctness repair used above.

The proof development used substantial assistance from large language models.  All mathematical claims in the public manuscript are intended to be presented in a self-contained human-checkable form.

\bibliographystyle{alpha}
\bibliography{references}

\end{document}
